\subsection{THE UNIMODULAR GROUP SL(n, A)}

Let \(\mathbf{M}_n(\mathbf{A})\) denote the ring of square matrices of order \(n\) over \(\mathbf{A}\) . Consider the mapping \(\det: \mathbf{M}_n(\mathbf{A}) \to \mathbf{A}\) . The group \(\mathbf{GL}(n, \mathbf{A})\) of invertible elements of \(\mathbf{M}_n(\mathbf{A})\) (isomorphic to the group of automorphisms of the A-module \(\mathbf{A}^n\) (II, § 10, no. 7)) is just the inverse image under this mapping of the multiplicative group \(\mathbf{A}^*\) of invertible elements of \(\mathbf{A}\) (no. 3, Proposition 5). Note on the other hand that the mapping \(\det: \mathbf{GL}(n, \mathbf{A}) \to \mathbf{A}^*\) is a group homomorphism (no. 3, Proposition 5).

The mapping \(\operatorname{det}:\mathbf{M}_n(\mathbf{A})\to \mathbf{A}\) is moreover surjective (and therefore so is the homomorphism \(\operatorname{det}:\mathbf{GL}(n,\mathbf{A})\to \mathbf{A}^*)\) ; since, for all \(\lambda \in \mathbf{A}\) ,

\[
\det (\operatorname{diag} (\lambda , 1, \dots , 1)) = \lambda
\]

by virtue of formula (32) of no. 6.

The kernel of the surjective homomorphism \(\operatorname{det}:\mathbf{GL}(n,\mathbf{A})\to \mathbf{A}^*\) is a normal subgroup of \(\mathbf{GL}(n,\mathbf{A})\) , which is composed of unimodular matrices; it is denoted by \(\mathbf{SL}_n(\mathbf{A})\) or \(\mathbf{SL}(n,\mathbf{A})\) and is often called the unimodular group or special linear group of square matrices of order \(n\) over \(\mathbf{A}\) .

In this no. we shall examine the case where A is a field. Recall that for \(1 \leqslant i \leqslant n\) , \(1 \leqslant j \leqslant n\) , \(E_{ij}\) denotes the square matrix of order \(n\) all of whose elements are zero except the one in the row of index \(i\) and the column of index \(j\) , which is equal to 1; with \(I_n\) denoting the unit matrix of order \(n\) , we write \(B_{ij}(\lambda) = I_n + \lambda E_{ij}\) for every ordered pair of distinct indices \(i, j\) and all \(\lambda \in A\) (II, § 10, no. 13).

PROPOSITION 17. Let K be a commutative field. The unimodular group \(\mathbf{SL}(n, \mathbf{K})\) is generated by the matrices \(B_{ij}(\lambda)\) with \(i \neq j\) and \(\lambda \in \mathbf{K}\) .

By II, § 10, no. 13, Proposition 14, we know that every matrix in \(\mathbf{GL}(n,\mathbf{K})\) is a product of matrices of the form \(B_{ij}(\lambda)\) and a matrix of the form \(\operatorname{diag}(1,1,\ldots,1,\alpha)\) with \(\alpha\in\mathbf{K}^{*}\) . Now it is immediate that \(\det(B_{ij}(\lambda))=1\) and \(\det(\operatorname{diag}(1,\ldots,1,\alpha))=\alpha\) (no. 6, Example 2); whence the proposition.

COROLLARY. The group \(\mathbf{SL}(n, \mathbf{K})\) is the group of commutators of \(\mathbf{GL}(n, \mathbf{K})\) , except in the case where \(n = 2\) and where \(\mathbf{K}\) is a field with 2 elements.

As \(\mathbf{SL}(n, \mathbf{K})\) is the kernel of the homomorphism det of \(\mathbf{GL}(n, \mathbf{K})\) to a commutative group \(\mathbf{K}^*\) , \(\mathbf{SL}(n, \mathbf{K})\) contains the commutator group \(\Gamma\) of \(\mathbf{GL}(n, \mathbf{K})\) (I, § 6, no. 2). To prove that \(\mathbf{SL}(n, \mathbf{K}) = \Gamma\) , it will suffice, by Proposition 17, to show that, for all \(\lambda \in \mathbf{K}^*\) , \(B_{ij}(\lambda)\) belongs to \(\Gamma\) . Now, \(B_{ij}(\lambda)\) is a conjugate of \(B_{ij}(1)\) in \(\mathbf{GL}(n, \mathbf{K})\) since \(B_{ij}(\lambda) = Q \cdot B_{ij}(1) \cdot Q^{-1}\) , where \(Q\) denotes the matrix with respect to the canonical basis ( \(e_i\) ) of the automorphism \(v\) of \(\mathbf{K}^n\) such that \(v(e_i) = \lambda e_i\) , \(v(e_k) = e_k\) for \(k \neq i\) . On the other hand, let \(u_{ij}\) (for \(i \neq j\) ) be the automorphism of \(\mathbf{K}^n\) such that \(u_{ij}(e_i) = -e_j\) , \(u_{ij}(e_j) = e_i\) , \(u_{ij}(e_k) = e_k\) for \(k \notin \{i, j\}\) , which belongs to \(\mathbf{SL}(n, \mathbf{K})\) ; then \(B_{ji}(\lambda) = U_{ij}B_{ij}(-\lambda)U_{ij}^{-1}\) , where \(U_{ij}\) is the matrix of \(u_{ij}\) with respect to the canonical basis. Similarly, if \(1 < i < j\) , then \(B_{1j}(\lambda) = U_{1i}B_{ij}(\lambda)U_{1i}^{-1}\) and finally, for \(2 < j\) , \(B_{12}(\lambda) = U_{2j}B_{1j}(\lambda)U_{2j}^{-1}\) . This proves that all the \(B_{ij}(\lambda)\) have the same image \(s\) in \(\mathbf{GL}(n, \mathbf{K}) / \Gamma\) and it remains to show that \(s\) is the identity element.

Suppose first that K contains an element \(\lambda\) distinct from 0 and 1; then \(1 = \lambda + (1 - \lambda)\) , the two terms on the right hand side being \(\neq 0\) ; the relation \(B_{12}(1) = B_{12}(\lambda)B_{12}(1 - \lambda)\) shows that \(s^{2} = s\) and hence s is the identity element.

Suppose now that \(n \geqslant 3\) . The product \(B_{21}(1)B_{31}(1)\) is the matrix of an automorphism \(u\) of \(\mathbf{K}^n\) such that \(u(e_1) = e_1 + e_2 + e_3\) , \(u(e_i) = e_i\) for \(i \neq 1\) . If \(S\) is the matrix of the automorphism \(u'\) of \(\mathbf{K}^n\) such that \(u'(e_2) = e_2 + e_3\) , \(u'(e_i) = e_i\) for \(i \neq 2\) , then \(S.B_{21}(1)B_{31}(1).S^{-1} = B_{21}(1)\) ; we also deduce that \(s^2 = s\) , which completes the proof.

Remarks. (1) \(\mathbf{GL}(2, \mathbf{F}_2) = \mathbf{SL}(2, \mathbf{F}_2)\) ; this is a solvable group of order 6, whose commutator group is of index 2 (II, § 10, Exercise 14).

\begin{enumerate}
\def\labelenumi{(\arabic{enumi})}
\setcounter{enumi}{1}
\tightlist
\item
  With the same notation as above, it can be proved as in I, § 5, no. 7, Proposition 9 that, for \(i < j\) , \(j - i > 1\) , \(u_{ij} = u_{j-1,j} u_{i,j-1} u_{j-1,j}^{-1}\) ; hence the group \(\mathbf{SL}(n, \mathbf{K})\) is generated by the matrices \(B_{12}(\lambda)\) and \(U_{i,i+1}\) for \(1 \leqslant i \leqslant n - 1\) .
\end{enumerate}

